%-------------------------
% Resume in Latex
% Author : Jake Gutierrez
% Based off of: https://github.com/sb2nov/resume
% License : MIT
%------------------------

\documentclass[letterpaper,11pt]{article}

\usepackage{latexsym}
\usepackage[empty]{fullpage}
\usepackage{titlesec}
\usepackage{marvosym}
\usepackage[usenames,dvipsnames]{color}
\usepackage{verbatim}
\usepackage{enumitem}
\usepackage[hidelinks]{hyperref}
\usepackage{fancyhdr}
\usepackage[english]{babel}
\usepackage{xeCJK}
\usepackage{fontspec}
\setmainfont{Arial}
\setCJKmainfont{Microsoft JhengHei}

% sans-serif
% \usepackage[sfdefault]{FiraSans}
% \usepackage[sfdefault]{roboto}
% \usepackage[sfdefault]{noto-sans}
% \usepackage[default]{sourcesanspro}

% serif
% \usepackage{CormorantGaramond}
% \usepackage{charter}


\pagestyle{fancy}
\fancyhf{} % clear all header and footer fields
\fancyfoot{}
\renewcommand{\headrulewidth}{0pt}
\renewcommand{\footrulewidth}{0pt}

% Adjust margins
\addtolength{\oddsidemargin}{-0.5in}
\addtolength{\evensidemargin}{-0.5in}
\addtolength{\textwidth}{1in}
\addtolength{\topmargin}{-0.7in}
\addtolength{\textheight}{1.4in}

\urlstyle{same}

\raggedbottom
\raggedright
\setlength{\tabcolsep}{0in}

% Sections formatting
\titleformat{\section}{
  \vspace{-4pt}\scshape\raggedright\large
}{}{0em}{}[\color{black}\titlerule \vspace{-5pt}]

% Ensure that generate pdf is machine readable/ATS parsable
% (Not needed for xelatex)


%-------------------------
% Custom commands
\newcommand{\resumeItem}[1]{
  \item\small{
    {#1 \vspace{-2pt}}
  }
}

\newcommand{\resumeSubheading}[4]{
  \vspace{-2pt}\item
    \begin{tabular*}{0.97\textwidth}[t]{l@{\extracolsep{\fill}}r}
      \textbf{#1} & #2 \\
      \textit{\small#3} & \textit{\small #4} \\
    \end{tabular*}\vspace{-7pt}
}

\newcommand{\resumeSubSubheading}[2]{
    \item
    \begin{tabular*}{0.97\textwidth}{l@{\extracolsep{\fill}}r}
      \textit{\small#1} & \textit{\small #2} \\
    \end{tabular*}\vspace{-7pt}
}

\newcommand{\resumeProjectHeading}[2]{
    \item
    \begin{tabular*}{0.97\textwidth}{l@{\extracolsep{\fill}}r}
      \small#1 & #2 \\
    \end{tabular*}\vspace{-7pt}
}

\newcommand{\resumeSubItem}[1]{\resumeItem{#1}\vspace{-4pt}}

\renewcommand\labelitemii{$\vcenter{\hbox{\tiny$\bullet$}}$}

\newcommand{\resumeSubHeadingListStart}{\begin{itemize}[leftmargin=0.15in, label={}]}
\newcommand{\resumeSubHeadingListEnd}{\end{itemize}}
\newcommand{\resumeItemListStart}{\begin{itemize}}
\newcommand{\resumeItemListEnd}{\end{itemize}\vspace{-5pt}}

%-------------------------------------------
%%%%%%  RESUME STARTS HERE  %%%%%%%%%%%%%%%%%%%%%%%%%%%%


\begin{document}

%----------HEADING----------

\begin{center}
    \textbf{\Huge \scshape 朱彥勳 (Patrick Zhu)} \\ \vspace{2pt}
    \small \href{mailto:Patrickzhu616@gmail.com}{\underline{Patrickzhu616@gmail.com}} $|$ 
    \href{https://github.com/Patrick-zhuyanxun}{\underline{github.com/Patrick-zhuyanxun}}
\end{center}


%-----------EDUCATION-----------
\section{學歷與證照 (Education \& Certifications)}
  \resumeSubHeadingListStart
    \resumeSubheading
      {國立陽明交通大學 (National Yang Ming Chiao Tung University)}{新竹, 台灣}
      {電控工程研究所 碩士}{2024年8月 -- 至今}

    \resumeSubheading
      {國立台北科技大學 (National Taipei University of Technology)}{台北, 台灣}
      {智慧自動化工程科 副學士 $|$ \textbf{GPA: 3.9/4.0}}{2018年8月 -- 2023年6月}
    \resumeSubheading
      {公務人員高等考試}{台灣}
      {電機工程技師}{2024年}
  \resumeSubHeadingListEnd


%-----------EXPERIENCE-----------
\section{工作經歷 (Experience)}
  \resumeSubHeadingListStart

    \resumeSubheading
      {和碩聯合科技 (Pegatron)}{2022年7月 -- 2022年8月}
      {約聘實習生 $|$ \emph{機構及工業設計中心}}{台北, 台灣}
      \resumeItemListStart
        \resumeItem{運用 Python 與 Open3D 函式庫開發 3D 點雲匹配系統 (GUI)，並參與量測用產線自走車之硬體機構設計}
      \resumeItemListEnd

    \resumeSubheading
      {精浚科技 (OME Technology Co., Ltd.)}{2021年7月 -- 2021年8月}
      {助理機構工程師 $|$ \emph{微奈米部門}}{台北, 台灣}
      \resumeItemListStart
        \resumeItem{負責壓電馬達定位平台與大型雷達設備之機構設計，並運用 SolidWorks 執行進階力學與運動模擬分析}
      \resumeItemListEnd

    \resumeSubheading
      {亞太美國學校 Pacific American School}{2025年}
      {機器人教練, Python教師}{新竹, 台灣}
      \resumeItemListStart
        \resumeItem{於亞太美國學校教導學生國際 VEX 機器人競賽}
        \resumeItem{教導學生關於機構設計、電路控制等方面，協助學生獲取參與世界賽資格}
      \resumeItemListEnd

  \resumeSubHeadingListEnd


%-----------COMPETITIONS-----------
\section{比賽與得獎經歷 (Competitions \& Awards)}
    \resumeSubHeadingListStart
      \resumeProjectHeading
          {\textbf{透過姿態辨識結合雙機械手臂進行瑕疵檢測} $|$ \emph{Python, 深度學習}}{2025年}
          \resumeItemListStart
            \resumeItem{榮獲 2025 經濟部 BestAI Awards AI 應用類學生組 \textbf{佳作}}
            \resumeItem{開發結合 PSO 微調的 6D 姿態預測系統，在 T-LESS 資料集上將 CosyPose 與 SC6D 的召回率分別提升 1.79\% 與 0.82\%，推論速度達 1 秒/張}
            \resumeItem{設計雙機械手臂協作與相機視角轉換與自動標註資料，提升標註效率與檢測視野與涵蓋率}
          \resumeItemListEnd
      \resumeProjectHeading
          {\textbf{永續能源智慧風扇控制系統} $|$ \emph{Python, IoT, LSTM}}{2022年}
          \resumeItemListStart
            \resumeItem{榮獲 2022 永續能源創意實作競賽 全國大專綠能創新組 \textbf{銀牌獎}}
            \resumeItem{處理溫濕度 IoT 感測數據，建立基於 LSTM 的時間序列預測模型動態調節 EC 風扇轉速，達成實質節能效率}
          \resumeItemListEnd
      \resumeProjectHeading
          {\textbf{第23屆、24屆 全國 TDK 盃機器人競賽}}{2019年 -- 2020年}
          \resumeItemListStart
            \resumeItem{榮獲 第24屆 \textbf{亞軍} 及 第23屆 \textbf{八強}}
          \resumeItemListEnd
    \resumeSubHeadingListEnd

%-----------PROJECTS-----------
\section{專案經歷 (Projects)}
    \resumeSubHeadingListStart
      \resumeProjectHeading
          {\textbf{語音識別機械手軟夾控制應用} $|$ \emph{Python, ROS 2, LLM, Whisper, PCL}}{2025年}
          \resumeItemListStart
            \resumeItem{整合 OpenAI Whisper、Qwen3B 與 YOLOE 建立語音指令物件辨識系統，成功率達 90\%}
            \resumeItem{利用 PCL 進行點雲轉換生成夾取姿態，並透過 ROS 2 MoveIt MTC 規劃手臂路徑，將整體系統反應時間（含語音分析 5 秒與路徑生成 5 秒）優化至 10 秒}
          \resumeItemListEnd
      \resumeProjectHeading
          {\textbf{智能垃圾桶 (Autonomous Trash Can)} $|$ \emph{ROS Noetic, Raspberry Pi, Arduino, MoveIt}}{2023年}
          \resumeItemListStart
            \resumeItem{與 MIT Media Lab 合作開發 (CSL UROP) APP 呼叫導航垃圾桶；負責核心機電整合與創新螺紋開蓋機構設計，串接 Raspberry Pi 與 Arduino 執行硬體控制}
            \resumeItem{運用 ROS MoveIt 與 RRT 演算法建構自主路徑規劃，並結合感測數據建立使用者垃圾量統計排行榜}
          \resumeItemListEnd

    \resumeSubHeadingListEnd

%
%-----------PROGRAMMING SKILLS-----------
\section{專業技能 (Technical Skills)}
 \begin{itemize}[leftmargin=0.15in, label={}]
    \small{\item{
     \textbf{程式語言}{: Python, C/C++, Linux Shell, AI agent } \\
     \textbf{框架與工具}{: ROS 1 / ROS 2, LeRobot, 樹莓派, IsaacSim, Pytorch, TensorFlow, MoveIt, Git, SolidWorks} \\
     \textbf{專業領域}{: 強化學習、影像處理、深度學習、機器人控制、機構設計、機電整合} \\
     \textbf{語言}{: 英文 (TOEIC: 800分)}
    }}
 \end{itemize}


%-------------------------------------------
\end{document}
